% ============================================================
% 分片：s10-gerbes  对应原文 PDF 第 119--124 页（印刷页 119--124）
% 负责章节：§10 Gerbes + 10.1--10.5（§10 标题起）
% 译者：trans-s07   审校状态：初稿
% ============================================================

% ==== 原文 PDF 第 119 页（印刷页 119）====
% [跨页接续: 原文第 119 页开头仍属 §9（由 s09-other 分片负责），本分片自该页 §10 标题起]

\section{gerbe}

\begin{keytakeaway}
\textbf{核心物理图像：}$G$ 与 ${}^{L}G$ 的镜像对称在一般纤维上给出对偶环面，但 Hitchin 纤维只是 torsor、没有典范基点，故对偶不能逐点陈述；正确的陈述是：$\mathcal{M}_H(SO_3)$ 的 B-模型须先被一个 $\mathbb{C}^\times$-gerbe $\mathcal{G}$ 扭转，才与 $\mathcal{M}_H(SL_2)$ 的 A-模型对偶。
\textbf{数学工具：}Hitchin 纤维化的截面（spin 结构与经典 opers）、典范丛的平方根、$\mathbb{C}^\times$/$\mathbb{Z}_2$-gerbe 与扭层、Pfaffian 线丛、轨形基本群与 $H^1(C,\mathbb{Z}_2)$、$\mathcal{G}$-扭 D-模与 Morita 等价双模。
\textbf{主要结论：}B-膜 $\mathcal{B}_+,\mathcal{B}_-$ 与 A-膜 $\boldsymbol{F}_1,\boldsymbol{F}_2$ 的配对没有典范选择，gerbe 扭转恰好吸收这一歧义；对应地，$SO_3$ 一侧的 A-模型须用 $\mathcal{G}$-扭 D-模，而 $SL_2$ 一侧因典范丛的平方根存在而不出现扭转。
\end{keytakeaway}

本节的目标是阐明 §2.2 中出现的一个微妙之点。若 $r$ 是 $\mathcal{M}_H(SO_3)$ 中对应于自同构群为 $\Gamma=\mathbb{Z}_2$ 的 $SO_3$ 局部系统的点，则有两个支撑在 $r$ 的膜，即 $\mathcal{B}_+$ 与 $\mathcal{B}_-$。相应的 $SL_2$ 的 Hitchin 纤维化的纤维是两个分支 $\boldsymbol{F}_1$ 与 $\boldsymbol{F}_2$ 的并，相应地在对偶 A-模型中有两个 A-膜。本文的核心主张是：A-膜 $\boldsymbol{F}_1$ 与 $\boldsymbol{F}_2$ 对偶于 B-膜 $\mathcal{B}_+$ 与 $\mathcal{B}_-$。但 $\boldsymbol{F}_1$ 与 $\boldsymbol{F}_2$ 中哪一个对偶于 $\mathcal{B}_+$、哪一个对偶于 $\mathcal{B}_-$？没有自然的方式决定；事实上，$\boldsymbol{F}_1$ 与 $\boldsymbol{F}_2$ 被对称群 $Q=\mathbb{Z}_2\times\mathbb{Z}_2$ 交换。相比之下，膜 $\mathcal{B}_+$ 与 $\mathcal{B}_-$ 并不等价：$\mathcal{B}_+$ 对应于 $\mathbb{Z}_2$ 的平凡表示，$\mathcal{B}_-$ 对应于非平凡表示。

\subsection{一个微妙之处}

% ==== 原文 PDF 第 120 页（印刷页 120）====

我们断言这一问题反映了 $\mathcal{M}_H(G)$ 与 $\mathcal{M}_H({}^{L}G)$ 的镜像对称中的一个微妙之处，它与内窥无关。让我们先问：Hitchin 纤维化是否有截面？\footnote{关于这类问题在更一般得多的语境中的研究，见 [DG]。}正如 [Hi1, Hi4] 所解释的，只要在 Riemann 面 $C$ 上选取一个 spin 结构，即 $C$ 的典范丛的平方根 $K^{1/2}$，就总能构造一个截面。令 $E=K^{-1/2}\oplus K^{1/2}$。尽管 $E$ 不稳定，但 Higgs 丛 $(E,\varphi)$ 仍可能是稳定的。这恰好发生在（在 $E$ 的一个唯一自同构意义下）
\begin{equation}\tag{10.1}
\varphi = \begin{pmatrix}0&1\\ w&0\end{pmatrix},
\end{equation}
其中 $w$ 是二次微分。此时 $\det\varphi=-w$，故对 $(E,\varphi)$ 在 Hitchin 纤维化下映到由二次微分 $-w$ 确定的底 $\boldsymbol{B}$ 中的点。由于 $\boldsymbol{B}$ 中每个点都来自唯一的这种形式的 $(E,\varphi)$，这给出 Hitchin 纤维化的一个截面。（这种形式的 Higgs 丛有时称为经典 opers，以反映它们与 [BD] 中 opers 的类比。）
% [Tutor Note: 推导补足。(10.1) 中的 $(E,\varphi)$ 是「经典 oper」的原型：$E=K^{-1/2}\oplus K^{1/2}$ 虽不稳定，但 $\varphi$ 把 $K^{1/2}$ 映入 $K^{-1/2}\otimes K_C=K^{1/2}$ 的分量恰好构成一个非零截面（即二次微分 $w$），使得 $(E,\varphi)$ 稳定且无 $\varphi$-不变子丛。Hitchin 映射把 $(E,\varphi)$ 送到 $\det\varphi=-w$ 所确定的一点，而 $w$ 取遍 $H^0(C,K_C^2)$ 的全体即得整个底 $\boldsymbol{B}$；因此截面与 spin 结构（即 $K^{1/2}$ 的选取）一一对应地构造出来。 来源：原书 p.120 上下文（A）。]

这样得到的截面并不完全典范，因为它依赖于 $K^{1/2}$ 的选取。然而，同样的构造（把 $E$ 换成 $H=\mathrm{ad}(E)$）对 $SO_3$ 有意义，而这里 $K^{1/2}$ 的选取无关紧要。于是对 $SO_3$，Hitchin 纤维化有自然的截面；但对 $SL_2$，截面的选取依赖于 $K^{1/2}$ 的选取。

这一区别实际上在 §2 的公式中可见。$SO_3$ 模空间在式 (3.39) 中描述；Hitchin 纤维化有一个由 $z=t=\infty$ 给出的自然截面。$SL_2$ 模空间在 (3.15) 中描述，而 Hitchin 纤维化没有自然截面。

对任意 $G$，重复这一构造：从经由 $\mathfrak{g}$ 的主 $\mathfrak{sl}_2$ 子代数的选取而与 $E$ 相联系的 $G$-丛出发。具有这种 $E$ 的 Higgs 丛 $(E,\varphi)$ 总给出 Hitchin 纤维化的一个截面。这个截面不依赖于 $K^{1/2}$ 的选取，当且仅当 $G$ 中对应于主 $\mathfrak{sl}_2$ 子代数的子群是 $SO_3$ 而非 $SL_2$。（例如，若 $G$ 是伴随型，则 Hitchin 纤维化有自然截面；$G=SL_{2n+1}$ 时同样如此。）

\subsection{镜像对称的一个疑难}

这些事实对 $\mathcal{M}_H(SO_3)$ 与 $\mathcal{M}_H(SL_2)$ 之间镜像对称的正确陈述造成了一个疑难。

设 ${}^{L}\boldsymbol{F}$ 与 $\boldsymbol{F}$ 是 $SO_3$ 与 $SL_2$ 的 Hitchin 纤维化的相应纤维。朴素地说，它们是对偶环面，意思是 ${}^{L}\boldsymbol{F}$ 参数化 $\boldsymbol{F}$ 上的平坦酉线丛，反之亦然。然而立刻有一个问题：紧复环面上的平坦酉线丛空间总有一个区分点，对应于平凡线丛。所以若 ${}^{L}\boldsymbol{F}$ 与 $\boldsymbol{F}$ 在这个意义下是对偶环面，则每一个都必须有一个区分点，与另一个上的平凡线丛相联系。这与下述事实矛盾：虽然 ${}^{L}\boldsymbol{F}$ 确实有区分点（它与上述 Hitchin 纤维化截面的交点），但 $\boldsymbol{F}$ 没有。事实上，[HT] 所证明的并不是 ${}^{L}\boldsymbol{F}$ 与 $\boldsymbol{F}$ 在这个朴素意义下对偶，而是阿贝尔簇 ${}^{L}\boldsymbol{F}$ 对偶于使 $\boldsymbol{F}$ 成为 torsor 的那个阿贝尔簇。
% [Tutor Note: 推导补足。这里被略去的一步是：紧复环面 $T$ 上的平坦酉线丛的模空间是 $\widehat{T}=\mathrm{Hom}(T,U_1)$，它有一个典范点（平凡线丛），故 $\widehat{T}$ 与 $T$ 之间的对偶逐点定义；但 Hitchin 纤维 $\boldsymbol{F}$ 仅是一个 $P_b$-torsor，没有基点，因此不能与 ${}^{L}\boldsymbol{F}$ 逐点对偶。正确的陈述是 $P_b$ 与 ${}^{L}\boldsymbol{F}$ 对偶（[HT] 的定理），而 $\boldsymbol{F}$ 的 torsor 结构正是「挑选基点」的自由度所在。下文引入的 gerbe 扭转恰好把这一自由度吸收进 B-模型的定义中。 来源：原书 p.120 上下文（A）。]

% ==== 原文 PDF 第 121 页（印刷页 121）====

对 $SO_3$ 与 $SL_2$ 理论之间的镜像对称，这意味着：$\mathcal{M}_H(SL_2)$ 的 A-模型并不完全对偶于 $\mathcal{M}_H(SO_3)$ 的 B-模型，而是对偶于这个 B-模型的一个略微扭转的版本。

复流形 $X$ 上的 B-膜是 $X$ 上的凝聚层（或相应导出范畴的对象，即这种层的复形，模去某个等价）。现在设 $\mathcal{G}$ 是 $X$ 上的 $\mathbb{C}^\times$ gerbe。对每个 $\mathcal{G}$，存在 B-膜范畴的 $\mathcal{G}$-扭版本；关于 $\mathcal{G}$-扭层的概念，见例如 [DJ] 第 2.3 节或 [Li] 定义 2.1.2.2。秩一的 $\mathcal{G}$-扭凝聚层是 $\mathcal{G}$ 的一个平凡化。$n$ 个平凡化的直和是秩 $n$ 的 $\mathcal{G}$-扭层的一个例子。

对我们当前的问题，需要在 $\mathcal{M}_H(SO_3)$ 上有一个平凡但非典范平凡的 $\mathbb{C}^\times$ gerbe。一般地，设 $X$ 是任意空间，$\mathcal{L}\to X$ 是复线丛。则存在一个典范定义的 gerbe $\mathcal{G}$，其（局部）平凡化是 $\mathcal{L}$ 的平方根。更确切地说，与开子集 $U\subset X$ 相联系的范畴的对象是二元组 $(\mathcal{M},\alpha)$，其中 $\mathcal{M}$ 是 $U$ 上的线丛，$\alpha$ 是 $\mathcal{M}^2$ 与 $\mathcal{L}$ 之间的同构。这个 gerbe $\mathcal{G}$ 整体平凡当且仅当 $\mathcal{L}$ 有整体平方根；并且它可以以唯一方式（差一个符号）平凡化，当且仅当 $\mathcal{L}$ 有唯一的整体平方根。$\mathcal{G}$ 是 $\mathbb{C}^\times$ gerbe，但它实际上通过嵌入 $\{\pm1\}\subset\mathbb{C}^\times$ 与一个 $\mathbb{Z}_2$ gerbe 相联系，因此有自然的平坦 gerbe 联络（用物理语言，它与 $X$ 上的平坦 B-场相联系）。

我们把这个构造应用于 $X=\mathcal{M}_H(G)$（$G$ 为某个约化 Lie 群）、$\mathcal{L}=K_X$（$X$ 的典范丛）的情形。考虑 $K_X$ 的平方根给出 $X$ 上的平坦 $\mathbb{C}^\times$ gerbe $\mathcal{G}$。这个 gerbe 实际上是平凡的，因为 $K_X$ 确实有整体平方根。这一点在 [BD] 第 4 节中有非常详细的解释，那里它出现的原因与当前的考虑有些类似。构造如下。给定 $C$ 上的任一 spin 结构 $\mathcal{S}$（即 $C$ 的典范丛的平方根）以及 $G$-丛 $E\to C$，考虑由 $\mathrm{ad}(E)$ 扭的、spin 结构 $\mathcal{S}$ 的 Dirac 算子的 Pfaffian 线丛 $\mathcal{L}_{\mathcal{S},E}$。当 $E$ 变动时，$\mathcal{L}_{\mathcal{S},E}$ 作为线丛 $\mathcal{L}_{\mathcal{S}}\to\mathcal{M}_H(G)$ 的纤维变动，后者是 $\mathcal{M}_H(G)$ 的典范丛的一个平方根。于是 gerbe $\mathcal{G}$ 对每个 spin 结构 $\mathcal{S}\to C$ 的选取都有自然的平凡化。对单连通的 $G$，$\mathcal{M}_H(G)$ 也是单连通的，因此这样得到的典范丛的平方根在不计同构意义下不依赖于 $\mathcal{S}$ 的选取。例如对 $G=SL_2$ 就是如此。在此情形，$\mathcal{G}$ 是典范平凡的。

一般地，若 $\mathcal{S}$ 与 $\mathcal{S}'$ 是 $C$ 上的两个 spin 结构，则 Pfaffian 构造给出 $\mathcal{M}_H(G)$ 的典范丛的两个平方根 $\mathcal{L}_{\mathcal{S}'}$ 与 $\mathcal{L}_{\mathcal{S}}$。它们必须相差与一个 2 阶线丛 $\mathcal{U}(\mathcal{S}',\mathcal{S})$ 的张量积：
\begin{equation}\tag{10.2}
\mathcal{L}_{\mathcal{S}'} = \mathcal{L}_{\mathcal{S}}\otimes\mathcal{U}(\mathcal{S}',\mathcal{S}).
\end{equation}
显然，对 $C$ 上三个 spin 结构 $\mathcal{S},\mathcal{S}',\mathcal{S}''$，我们有
\begin{equation}\tag{10.3}
\mathcal{U}(\mathcal{S}'',\mathcal{S}) = \mathcal{U}(\mathcal{S}'',\mathcal{S}')\otimes\mathcal{U}(\mathcal{S}',\mathcal{S}).
\end{equation}
一个特别简单的例子是 $G=SO_3$。$H^1(C,\mathbb{Z}_2)$ 与轨形基本群 $\pi_1(\mathcal{M}_H(SO_3))$ 之间有自然的同构。于是有自然的映射，把 2 阶线丛 $\mathcal{V}\to C$ 映到 2 阶轨形线丛 $\mathcal{T}(\mathcal{V})\to\mathcal{M}_H(SO_3)$。若 $\mathcal{S}',\mathcal{S}\to C$ 是两个 spin 结构，则 $\mathcal{S}'=\mathcal{S}\otimes\mathcal{V}$ 对某个 2 阶线丛 $\mathcal{V}$ 成立，并且
\begin{equation}\tag{10.4}
\mathcal{U}(\mathcal{S}',\mathcal{S}) = \mathcal{T}(\mathcal{V}),
\end{equation}
这是与式 (10.3) 显然相容的陈述。
% [Tutor Note: 物理图像。gerbe 的「平坦性」来自它经由 $\{\pm1\}\subset\mathbb{C}^\times$ 与 $\mathbb{Z}_2$-gerbe 相联系：$\mathbb{Z}_2$ 是离散群，故对应的 gerbe 联络是平坦的（即 B-场是平坦的）。不同 spin 结构给出的平凡化相差 2 阶线丛 $\mathcal{U}(\mathcal{S}',\mathcal{S})$，其类正是 $H^1(C,\mathbb{Z}_2)$ 中的元素经 $\mathcal{T}$ 的像；这解释了为什么镜像对称中出现的扭转恰好是 $\mathbb{Z}_2$ 型的、且与 spin 结构的选取无关（只差一个可吸收的 2 阶线丛）。 来源：原书 p.121 上下文（A）。]

% ==== 原文 PDF 第 122 页（印刷页 122）====

$\mathcal{M}_H(SO_3)$ 与 $\mathcal{M}_H(SL_2)$ 之间镜像对称的精确陈述如下。$SO_3$ 的 $\mathcal{G}$-扭 B-模型对偶于 $SL_2$ 的 A-模型。反之，$\mathcal{M}_H(SO_3)$ 的 $\mathcal{G}$-扭 A-模型对偶于 $\mathcal{M}_H(SL_2)$ 的 B-模型。一般地，A-模型可以被平坦复 gerbe 扭转，例如 $\mathcal{G}$。

类似的 gerbe 扭转也应当在 Beilinson 与 Drinfeld 提出的几何 Langlands 对应的非阿贝尔 Fourier--Mukai 变换表述中实现。

\subsection{应用}

现在让我们重新考虑本节开头的问题：B-膜 $\mathcal{B}_+$ 与 $\mathcal{B}_-$ 和 A-膜 $\boldsymbol{F}_1$ 与 $\boldsymbol{F}_2$ 之间的对偶。哪一个 A-膜对应于 $\mathcal{B}_+$？

首先，问题之所以存在，只是因为 $SL_2$ 的 Hitchin 纤维化没有自然截面。给定这样一个截面 $s$，我们就能挑出特殊 Hitchin 纤维的区分分支 $\boldsymbol{F}_1$ 或 $\boldsymbol{F}_2$，即与 $s$ 相交的那一个。问题的解决在于：对偶并不涉及 $\mathcal{M}_H(SO_3)$ 的普通 B-模型，而涉及 $\mathcal{G}$-扭 B-模型。在普通 B-模型中，$\mathcal{B}_+$ 与 $\mathcal{B}_-$ 之间有一个区分出来的，即 $SO_3$ 局部系统的自同构群 $\mathbb{Z}_2$ 平凡地作用的那个。但在 $\mathcal{G}$-扭 B-模型中，情况不同。虽然两个 B-膜在 $\Gamma=\mathbb{Z}_2$ 下以相反方式变换，但要弄清哪一个平凡地变换、哪一个非平凡地变换，我们必须先平凡化 $\mathcal{G}$。

然而，$\mathcal{G}$ 没有自然的平凡化；相反，它有一族依赖于 $C$ 上 spin 结构选取的平凡化。不同的平凡化会给出关于两个 B-膜中哪一个在 $\Gamma$ 下不变、哪一个不是的不同解释。

事实上，这些平凡化中的两个相差与线丛 $\mathcal{T}(\mathcal{V})$ 之一的张量积，其中 $\mathcal{V}\in H^1(C,\mathbb{Z}_2)$。若 $r\in\mathcal{M}_H(SO_3)$ 是轨形奇点之一，对称群为 $\Gamma=\mathbb{Z}_2$，则对适当的 $\mathcal{V}$，$\Gamma$ 的非平凡元素作用在 $\mathcal{T}(\mathcal{V})$ 的纤维上为乘以 $-1$。在这种情形，哪一个 B-膜是 $\Gamma$-不变的、哪一个不是，这一选择会被与 $\mathcal{T}(\mathcal{V})$ 的张量积反转。

\subsection{对偶对称群}

为了更多的理解，我们应当描述我们所利用的某些事实的另一种解释。

群 $Q=H^1(C,\mathbb{Z}_2)$ 以 §3.5 中熟悉的方式作用在 $\mathcal{M}_H(SL_2)$ 上：$Q$ 的元素对应于 2 阶线丛 $\mathcal{V}\to C$，它通过 $E\to E\otimes\mathcal{V}$ 作用在 Higgs 丛 $(E,\varphi)$ 上。这给出 $Q$ 在 $\mathcal{M}_H(SL_2)$ 上的几何作用，保持其超凯勒结构，因而给出 $Q$ 在 $\mathcal{M}_H(SL_2)$ 的 A-模型与 B-模型上的作用。

对偶地，同构的群必须作用\footnote{另见 [KW] 第 7.2 节中的另一种解释。}在对偶模空间 $\mathcal{M}_H(SO_3)$ 的 B-模型与 A-模型上。关键在于我们上面已经利用过的：

% ==== 原文 PDF 第 123 页（印刷页 123）====

从 2 阶线丛 $\mathcal{V}\to C$ 到 2 阶线丛 $\mathcal{T}(\mathcal{V})\to\mathcal{M}_H(SO_3)$ 的自然对应。由于 $\mathcal{T}(\mathcal{V})$ 是凝聚层，与 $\mathcal{T}(\mathcal{V})$ 的张量积作为 B-模型的对称有意义；由于它是平坦线丛，它作为 A-模型的对称也有意义。

$\mathcal{M}_H(SL_2)$ 与 $\mathcal{M}_H(SO_3)$ 之间的对偶，把 $Q$ 在 $\mathcal{M}_H(SL_2)$ 的 A-与 B-模型上由其在空间上的几何作用给出的作用，与 $Q$ 通过张量积在 $\mathcal{M}_H(SO_3)$ 的 $\mathcal{G}$-扭 B-与 A-模型上的作用互相交换。

$Q$ 出于几何理由交换支撑在 $\mathcal{M}_H(SL_2)$ 特殊 Hitchin 纤维上的两个 A-膜。它交换支撑在轨形奇点上的两个 B-膜，是因为 $\Gamma=\mathbb{Z}_2$ 的非平凡元素作用在某些线丛 $\mathcal{T}(\mathcal{V})$ 的相应纤维上为 $-1$。

\subsection{与几何 Langlands 通常陈述的关系}

在本文中，为了探究内窥，我们主要把 $SO_3$ 的 $\mathcal{G}$-扭 B-模型与 $SL_2$ 的 A-模型相比较。然而，把 $SL_2$ 的 B-模型与 $SO_3$ 的 $\mathcal{G}$-扭 A-模型相比较也很有意思。$\mathcal{G}$-扭转在那个语境中起什么作用？

[KW] 第 11 节表明，对任意 $G$，$\mathcal{M}_H(G)$ 上的 A-膜等价于 $\mathcal{M}(G)$（$G$-丛的模空间）上的扭 D-模。这里\footnote{为简单起见，我们只考虑非分歧情形。分歧会导致进一步的扭转，但不影响我们的主要主张。}，扭 D-模是某个代数层的模层，该代数层我们记作 $\CD^*$，即作用在 $\mathcal{M}$ 的典范丛 $K_{\mathcal{M}}$ 的一个平方根上的微分算子。代数层 $\CD^*$ 不依赖于 $K_{\mathcal{M}}$ 平方根的整体选取（甚至不依赖于这种平方根的整体存在性，尽管它事实上存在）。

这与几何 Langlands 对偶的通常陈述之间有轻微的紧张：对偶的右端本应涉及普通 D-模，而不是扭 D-模。

现在，若存在 $K_{\mathcal{M}}$ 的典范整体平方根，则这一区别无关紧要。给定这样的线丛，我们可以考虑把 $\mathcal{O}_{\mathcal{M}}$ 映到 $K_{\mathcal{M}}^{1/2}$ 的微分算子，这种算子的层是（代数层的）对 $(\CD,\CD^*)$——即普通微分算子与扭微分算子——的「双模」。这个双模是「Morita 等价双模」，它建立普通 D-模范畴与扭 D-模范畴之间的等价。

对 $SL_2$，存在这样的 $K_{\mathcal{M}}^{1/2}$ 的典范选取；但对 $SO_3$ 没有。于是 $\mathcal{M}(SO_3)$ 上的 $\CD^*$-模（例如我们从 $\mathcal{M}_H(SO_3)$ 的 A-膜得到的）并不典范地等同于 $\mathcal{M}(SO_3)$ 上的普通 D-模（例如我们在几何 Langlands 纲领中所期望的）。

调和这两种观点的是：在 S-对偶中出现的 $\mathcal{M}_H(SO_3)$ 上的 A-模型是 $\mathcal{G}$-扭的。普通的 A-膜映到 $\mathcal{M}(SO_3)$ 上的 $\CD^*$-模，而 $\mathcal{G}$-扭 A-膜映到同一空间上的 $\mathcal{G}$-扭 $\CD^*$-模。但 $\mathcal{M}(SO_3)$ 上的 $\mathcal{G}$-扭 $\CD^*$-模典范地映到同一空间上的普通 D-模。原因在于：尽管平方根 $K_{\mathcal{M}}^{1/2}$ 作为线丛不典范地存在，它却作为 gerbe $\mathcal{G}$ 的一个平凡化而典范地——同义反复地——存在。
% [Tutor Note: 易错预警。这一段把「扭转」的两个不同层次混在一起，阅读时须区分：(1) B-模型的 $\mathcal{G}$-扭来自典范丛平方根的整体不存在性（§10.2）；(2) 从 A-膜到 D-模的过渡中，$\CD^*$ 是作用在 $K_{\mathcal{M}}^{1/2}$ 上的微分算子层，它作为代数层并不依赖平方根的整体选取。关键是「$K_{\mathcal{M}}^{1/2}$ 作为 $\mathcal{G}$ 的平凡化典范存在」：这使 $\mathcal{G}$-扭 $\CD^*$-模与普通 D-模等价，从而 $SO_3$ 一侧的扭 A-膜仍给出几何 Langlands 通常意义下的 D-模。$SL_2$ 情形因为 $K_{\mathcal{M}}^{1/2}$ 本身有典范选取，扭转自动消失。 来源：原书 p.123 上下文（A）。]

% ==== 原文 PDF 第 124 页（印刷页 124）====

我们已对作为本文主要例子的对偶群 $SO_3$ 与 $SL_2$ 描述了这一点。更一般地，对任意约化群 $G$，人们定义典范丛平方根的 gerbe $\mathcal{G}$。给定任一对偶对 ${}^{L}G$ 与 $G$，底层的规范场论对偶是 ${}^{L}G$ 的 $\mathcal{G}$-扭 B-模型与 $G$ 的 $\mathcal{G}$-扭 A-模型之间的同构（以及反之）。在我们的例子中，我们只在一边看到扭转，仅仅是因为 gerbe $\mathcal{G}$ 对 $SL_2$ 是典范平凡的。

% [以下为 §11 附录标题起的内容，属 s11-appendix 分片]
